\documentclass{beamer}

\usepackage{amsmath,amsfonts,amssymb,bm}
\usepackage{quantikz}
\usepackage{braket}
\usepackage{physics}
\usepackage{tikz}
\usetikzlibrary{arrows.meta,positioning,calc}

\usetheme{Madrid}

\title{The Quantum Fourier Transform}
\subtitle{Two-Qubit Derivation, Binary Fractions, and Phase Estimation}
\author{}
\date{}

\begin{document}

\frame{\titlepage}

%================================================
\section{Motivation}
%================================================

\begin{frame}{Why the Quantum Fourier Transform?}
The quantum Fourier transform (QFT) is one of the central subroutines in quantum computing.

\medskip

It appears in:
\begin{itemize}
\item quantum phase estimation,
\item Shor's factoring algorithm,
\item hidden subgroup algorithms,
\item period finding and spectral analysis.
\end{itemize}

\medskip

Mathematically, it is the unitary map
\[
\mathrm{QFT}_N \ket{x}
=
\frac{1}{\sqrt{N}}
\sum_{k=0}^{N-1}
e^{2\pi i xk/N}\ket{k},
\qquad N=2^n.
\]
\end{frame}

%------------------------------------------------

\begin{frame}{Main Goals of These Notes}
We will:
\begin{enumerate}
\item derive the two-qubit QFT step by step,
\item explain why the output appears in reversed order,
\item generalize the circuit to \(n\) qubits,
\item show the role of binary fractions,
\item connect the QFT directly to quantum phase estimation.
\end{enumerate}
\end{frame}

%================================================
\section{Definition of the QFT}
%================================================

\begin{frame}{Definition}
For \(N=2^n\), the QFT acts on the computational basis as
\[
\mathrm{QFT}_N \ket{x}
=
\frac{1}{\sqrt{N}}
\sum_{k=0}^{N-1}
\omega_N^{xk}\ket{k},
\qquad
\omega_N=e^{2\pi i/N}.
\]

If
\[
x = x_{n-1}2^{n-1}+x_{n-2}2^{n-2}+\cdots+x_1 2 + x_0,
\]
with \(x_j\in\{0,1\}\), then \(\ket{x}\equiv \ket{x_{n-1}\cdots x_0}\).
\end{frame}

%------------------------------------------------

\begin{frame}{Binary Fractions}
A compact way to write QFT output phases uses binary fractions.

\medskip

For bits \(x_1,x_0\), define
\[
0.x_1x_0 = \frac{x_1}{2}+\frac{x_0}{4}.
\]

More generally,
\[
0.x_{m}x_{m-1}\cdots x_0
=
\frac{x_m}{2}+\frac{x_{m-1}}{2^2}+\cdots+\frac{x_0}{2^{m+1}}.
\]

These binary fractions will appear naturally in the tensor-product form of the QFT.
\end{frame}

%================================================
\section{Two-Qubit QFT: Algebraic Derivation}
%================================================

\begin{frame}{Two-Qubit Input State}
Let
\[
\ket{x}=\ket{x_1x_0},
\qquad x=2x_1+x_0,
\]
where \(x_1\) is the most significant bit and \(x_0\) the least significant bit.

For two qubits, \(N=4\), so
\[
\mathrm{QFT}_4 \ket{x}
=
\frac{1}{2}
\sum_{k=0}^{3}
e^{2\pi i xk/4}\ket{k}.
\]
\end{frame}

%------------------------------------------------

\begin{frame}{Expand the Sum Explicitly}
Write \(k=2k_1+k_0\), with \(k_1,k_0\in\{0,1\}\). Then
\[
\mathrm{QFT}_4 \ket{x_1x_0}
=
\frac{1}{2}
\sum_{k_1,k_0\in\{0,1\}}
e^{2\pi i (2x_1+x_0)(2k_1+k_0)/4}
\ket{k_1k_0}.
\]

Expand the exponent:
\[
\frac{(2x_1+x_0)(2k_1+k_0)}{4}
=
x_1k_1 + \frac{x_1k_0}{2} + \frac{x_0k_1}{2} + \frac{x_0k_0}{4}.
\]

Since \(e^{2\pi i x_1k_1}=1\) for integers \(x_1k_1\), we obtain
\[
e^{2\pi i xk/4}
=
e^{2\pi i \left( \frac{x_1k_0}{2} + \frac{x_0k_1}{2} + \frac{x_0k_0}{4}\right)}.
\]
\end{frame}

%------------------------------------------------

\begin{frame}{Factorization into a Tensor Product}
Now separate the dependence on \(k_1\) and \(k_0\):
\[
e^{2\pi i xk/4}
=
e^{2\pi i k_1 \left(\frac{x_0}{2}\right)}
\,
e^{2\pi i k_0 \left(\frac{x_1}{2}+\frac{x_0}{4}\right)}.
\]

Hence
\[
\mathrm{QFT}_4 \ket{x_1x_0}
=
\frac{1}{2}
\sum_{k_1,k_0}
e^{2\pi i k_1(0.x_0)}
e^{2\pi i k_0(0.x_1x_0)}
\ket{k_1k_0}.
\]

This factorizes as
\[
\mathrm{QFT}_4 \ket{x_1x_0}
=
\frac{1}{\sqrt{2}}
\left(\ket{0}+e^{2\pi i\,0.x_0}\ket{1}\right)
\otimes
\frac{1}{\sqrt{2}}
\left(\ket{0}+e^{2\pi i\,0.x_1x_0}\ket{1}\right).
\]
\end{frame}

%------------------------------------------------

\begin{frame}{Interpretation of the Two Output Qubits}
We have found
\[
\mathrm{QFT}_4 \ket{x_1x_0}
=
\frac{1}{\sqrt{2}}
\left(\ket{0}+e^{2\pi i\,0.x_0}\ket{1}\right)
\otimes
\frac{1}{\sqrt{2}}
\left(\ket{0}+e^{2\pi i\,0.x_1x_0}\ket{1}\right).
\]

Key point:
\begin{itemize}
\item the first factor depends only on \(x_0\),
\item the second factor depends on \(x_1x_0\).
\end{itemize}

So the least significant phase appears first. This is why the circuit output is naturally in reversed order.
\end{frame}

%================================================
\section{Two-Qubit Circuit}
%================================================

\begin{frame}{Useful Elementary Gates}
We use:
\[
H\ket{0}=\frac{\ket{0}+\ket{1}}{\sqrt{2}},
\qquad
H\ket{1}=\frac{\ket{0}-\ket{1}}{\sqrt{2}},
\]
and the controlled phase gate
\[
R_2=
\begin{pmatrix}
1&0\\
0&e^{2\pi i/2^2}
\end{pmatrix}
=
\begin{pmatrix}
1&0\\
0&e^{i\pi/2}
\end{pmatrix}.
\]

For two qubits, the QFT uses one Hadamard on each qubit, one controlled-\(R_2\), and one final SWAP.
\end{frame}

%------------------------------------------------

\begin{frame}{Two-Qubit QFT Circuit}
\centering
\begin{quantikz}[row sep=0.35cm, column sep=0.5cm]
\lstick{$\ket{x_1}$} & \gate{H} & \ctrl{1}   & \qw      & \swap{1}   & \qw \\
\lstick{$\ket{x_0}$} & \qw      & \gate{R_2} & \gate{H} & \targX{}{} & \qw
\end{quantikz}

\medskip

Reading from left to right:
\begin{enumerate}
\item apply \(H\) to the most significant qubit,
\item apply a controlled \(R_2\),
\item apply \(H\) to the least significant qubit,
\item swap the two output qubits.
\end{enumerate}
\end{frame}

%------------------------------------------------

\begin{frame}{Step 1: Apply \(H\) to the First Qubit}
Start with
\[
\ket{x_1x_0}.
\]

Applying \(H\) to the first qubit gives
\[
\ket{x_1x_0}
\mapsto
\frac{1}{\sqrt{2}}
\left(\ket{0}+e^{i\pi x_1}\ket{1}\right)\otimes \ket{x_0}.
\]

Since \(e^{i\pi x_1}=(-1)^{x_1}\), this is
\[
\frac{1}{\sqrt{2}}
\left(\ket{0}+e^{2\pi i \, x_1/2}\ket{1}\right)\otimes \ket{x_0}.
\]
\end{frame}

%------------------------------------------------

\begin{frame}{Step 2: Controlled Phase \(R_2\)}
The controlled-\(R_2\) adds a phase \(e^{i\pi/2}\) only when both qubits are in \(\ket{1}\).

Thus the amplitude of \(\ket{1}\otimes\ket{x_0}\) picks up an extra factor
\[
e^{2\pi i\, x_0/4}.
\]

So the state becomes
\[
\frac{1}{\sqrt{2}}
\left(
\ket{0}
+
e^{2\pi i \left(\frac{x_1}{2}+\frac{x_0}{4}\right)}
\ket{1}
\right)\otimes \ket{x_0}.
\]

That is
\[
\frac{1}{\sqrt{2}}
\left(
\ket{0}
+
e^{2\pi i\, 0.x_1x_0}\ket{1}
\right)\otimes \ket{x_0}.
\]
\end{frame}

%------------------------------------------------

\begin{frame}{Step 3: Hadamard on the Second Qubit}
Now apply \(H\) to the second qubit:
\[
\ket{x_0}
\mapsto
\frac{1}{\sqrt{2}}
\left(\ket{0}+e^{2\pi i\,x_0/2}\ket{1}\right).
\]

Hence the full state becomes
\[
\frac{1}{\sqrt{2}}
\left(
\ket{0}
+
e^{2\pi i\,0.x_1x_0}\ket{1}
\right)
\otimes
\frac{1}{\sqrt{2}}
\left(
\ket{0}
+
e^{2\pi i\,0.x_0}\ket{1}
\right).
\]
\end{frame}

%------------------------------------------------

\begin{frame}{Step 4: Why the SWAP is Needed}
At this stage, the circuit has produced
\[
\frac{1}{\sqrt{2}}
\left(
\ket{0}
+
e^{2\pi i\,0.x_1x_0}\ket{1}
\right)
\otimes
\frac{1}{\sqrt{2}}
\left(
\ket{0}
+
e^{2\pi i\,0.x_0}\ket{1}
\right),
\]
which is the correct QFT state but with the qubits in reversed order compared with the standard algebraic expression.

The final SWAP changes
\[
\ket{q_1}\ket{q_0}\longrightarrow \ket{q_0}\ket{q_1},
\]
restoring the conventional most-significant-to-least-significant ordering.
\end{frame}

%================================================
\section{Generalization to \(n\) Qubits}
%================================================

\begin{frame}{General Tensor-Product Form}
For
\[
\ket{x}=\ket{x_{n-1}x_{n-2}\cdots x_1x_0},
\]
the QFT can be written as
\[
\mathrm{QFT}_{2^n}\ket{x}
=
\frac{1}{2^{n/2}}
\bigotimes_{j=0}^{n-1}
\left(
\ket{0}
+
e^{2\pi i\,0.x_jx_{j-1}\cdots x_0}\ket{1}
\right),
\]
up to reversal of qubit order.

Equivalently, in standard order,
\[
\mathrm{QFT}_{2^n}\ket{x}
=
\bigotimes_{m=0}^{n-1}
\frac{1}{\sqrt{2}}
\left(
\ket{0}
+
e^{2\pi i\,0.x_{n-1-m}\cdots x_0}\ket{1}
\right),
\]
followed by a bit-reversal.
\end{frame}

%------------------------------------------------

\begin{frame}{Meaning of the Binary Fractions}
The phases in the \(n\)-qubit QFT are:
\[
0.x_0 = \frac{x_0}{2},
\]
\[
0.x_1x_0 = \frac{x_1}{2}+\frac{x_0}{4},
\]
\[
0.x_2x_1x_0 = \frac{x_2}{2}+\frac{x_1}{4}+\frac{x_0}{8},
\]
and so on.

Thus each output qubit stores a phase that depends on a progressively longer suffix of the input bit string.
\end{frame}

%------------------------------------------------

\begin{frame}{General \(n\)-Qubit QFT Circuit}
For \(n\) qubits, the QFT circuit is built recursively:
\begin{itemize}
\item apply \(H\) to the first qubit,
\item apply controlled-\(R_2, R_3,\dots,R_n\) from lower qubits,
\item move to the next qubit,
\item finish with pairwise swaps.
\end{itemize}

The controlled phase gates are
\[
R_k=
\begin{pmatrix}
1 & 0 \\
0 & e^{2\pi i/2^k}
\end{pmatrix}.
\]
\end{frame}

%------------------------------------------------

\begin{frame}{Schematic \(n\)-Qubit Circuit}
\centering
\begin{quantikz}[row sep=0.25cm, column sep=0.3cm]
\lstick{$\ket{x_{n-1}}$} & \gate{H} & \ctrl{1}   & \ctrl{2}   & \qw      & \cdots & \ctrl{3} & \qw \\
\lstick{$\ket{x_{n-2}}$} & \qw      & \gate{R_2} & \qw        & \gate{H} & \cdots & \ctrl{2} & \qw \\
\lstick{$\ket{x_{n-3}}$} & \qw      & \qw        & \gate{R_3} & \qw      & \cdots & \ctrl{1} & \qw \\
\vdots                 &          &            &            &          & \ddots &            &     \\
\lstick{$\ket{x_0}$}     & \qw      & \qw        & \qw        & \qw      & \cdots & \gate{R_n} & \gate{H} & \qw
\end{quantikz}

followed by swaps
\[
(1\leftrightarrow n),\quad
(2\leftrightarrow n-1),\quad \dots
\]
to reverse the order of the output qubits.
\end{frame}

%------------------------------------------------

\begin{frame}{Gate Count}
The exact \(n\)-qubit QFT uses:
\begin{itemize}
\item \(n\) Hadamard gates,
\item \(\frac{n(n-1)}{2}\) controlled phase gates,
\item \(\lfloor n/2\rfloor\) swaps.
\end{itemize}

So the exact gate complexity is
\[
\mathcal{O}(n^2).
\]

If very small controlled phases are neglected, one gets the approximate QFT, which can reduce the gate count significantly.
\end{frame}

%================================================
\section{Connection to Phase Estimation}
%================================================

\begin{frame}{Why the QFT Appears in Phase Estimation}
Suppose
\[
U\ket{u}=e^{2\pi i \phi}\ket{u},
\qquad 0\le \phi<1.
\]

Quantum phase estimation aims to determine the binary digits of \(\phi\).

The algorithm creates a superposition of time labels:
\[
\frac{1}{2^{m/2}}\sum_{k=0}^{2^m-1}\ket{k}\ket{u},
\]
and controlled powers of \(U\) convert this into
\[
\frac{1}{2^{m/2}}\sum_{k=0}^{2^m-1} e^{2\pi i k\phi}\ket{k}\ket{u}.
\]

This is a Fourier-type state in the control register.
\end{frame}

%------------------------------------------------

\begin{frame}{Inverse QFT Extracts the Phase Bits}
If
\[
\phi = 0.\phi_1\phi_2\cdots \phi_m
\]
has an exact \(m\)-bit binary expansion, then the inverse QFT maps
\[
\frac{1}{2^{m/2}}
\sum_{k=0}^{2^m-1} e^{2\pi i k\phi}\ket{k}
\longrightarrow
\ket{\phi_1\phi_2\cdots \phi_m}.
\]

So the QFT and inverse QFT are the natural transformations that convert between:
\begin{itemize}
\item phase-encoded superpositions,
\item binary computational-basis strings.
\end{itemize}
\end{frame}

%------------------------------------------------

\begin{frame}{Binary Fractions in Phase Estimation}
The measured output bits are precisely the binary fraction digits:
\[
\phi = 0.\phi_1\phi_2\cdots \phi_m
=
\frac{\phi_1}{2}+\frac{\phi_2}{4}+\cdots+\frac{\phi_m}{2^m}.
\]

This is exactly the same binary-fraction structure that appears in the QFT tensor-product decomposition.
\end{frame}

%------------------------------------------------

\begin{frame}{Two-Qubit Phase Estimation Example}
With two control qubits, suppose
\[
\phi = 0.\phi_1\phi_0
=
\frac{\phi_1}{2}+\frac{\phi_0}{4}.
\]

Then after the controlled powers of \(U\), the control register is in the state
\[
\frac{1}{2}
\sum_{k=0}^{3} e^{2\pi i k\phi}\ket{k}.
\]

Applying \(\mathrm{QFT}_4^{-1}\) yields
\[
\ket{\phi_1\phi_0},
\]
assuming the phase is exactly representable with two binary digits.
\end{frame}

%------------------------------------------------

\begin{frame}{Conceptual Link}
The QFT reorganizes information in a way that is natural for phases.

\medskip

In phase estimation:
\begin{itemize}
\item controlled powers of \(U\) write the eigenphase into amplitudes,
\item the inverse QFT converts those phase patterns into a binary number.
\end{itemize}

So the deep link is:
\[
\text{QFT} \leftrightarrow \text{binary fractions} \leftrightarrow \text{phase readout}.
\]
\end{frame}

%================================================
\section{Geometric and Physical Interpretation}
%================================================

\begin{frame}{Complex-Phase View}
Each output qubit of the QFT is a qubit on the equator of the Bloch sphere:
\[
\frac{1}{\sqrt{2}}
\left(
\ket{0}+e^{i\theta}\ket{1}
\right).
\]

The angle \(\theta\) is determined by a binary fraction of the input bits.

Thus the QFT converts a computational basis state into a tensor product of phase states.
\end{frame}

%------------------------------------------------

\begin{frame}{Why the Output is Reversed}
The circuit builds the least significant phase first:
\[
0.x_0,\quad 0.x_1x_0,\quad 0.x_2x_1x_0,\dots
\]

Therefore the qubits emerge in little-endian order.

The final swaps simply restore the conventional ordering of the qubits. In many algorithms, these swaps can be omitted if one keeps track of the reversed labeling classically.
\end{frame}

%================================================
\section{Summary}
%================================================

\begin{frame}{Summary}
\begin{itemize}
\item For two qubits,
\[
\mathrm{QFT}_4\ket{x_1x_0}
=
\frac{1}{\sqrt{2}}
(\ket{0}+e^{2\pi i\,0.x_0}\ket{1})
\otimes
\frac{1}{\sqrt{2}}
(\ket{0}+e^{2\pi i\,0.x_1x_0}\ket{1}),
\]
up to output reversal.
\item The circuit uses Hadamards, controlled phase gates, and a final SWAP.
\item For \(n\) qubits, the QFT produces one output qubit per binary fraction suffix.
\item This binary-fraction structure is exactly what makes the QFT central in phase estimation.
\end{itemize}
\end{frame}

\end{document}
